\chapter{相对论电动力学}
\label{chap:relativistic electromagnetic}

%三维一般波动方程
%\laplacian u-\frac1{a^2}\pv[2]ut=-f(\bm r,t),
%中的$a$是波速，依惯性系而不同。

\paragraph{Michelson--Morley实验}

人们最开始设想光需要在一种称为以太(ether)的介质中传播，光速是相对于以太的速度。
为了验证以太是否存在以及光速是否是常数，1887年Michelson和Morley设计了如下所示的光学干涉装置：
\begin{notecaptionblock}
	\tikzfigure{michelson-morley-experiment}
	\captionof{figure}{Michelson--Morley 光学干涉仪。两束光被端镜 $\mathrm M_1,\mathrm M_2$ 反射后回到分束器 BS（虚线），在探测端 $\mathrm D$ 观察干涉。短斜线标出镜背。}
	\label{fig:Michelson-Morley}
\end{notecaptionblock}
图中最左侧相干光源发出的光在经过分束器(beam splitter, BS)中心$O$后，其中一部分光被透射后抵达镜子M$_1$然后被反射回$O$点，另一部分光被反射后抵达镜子M$_2$然后被反射回$O$点，并与从M$_1$反射回来的光发生干涉。假设整个装置在以太中以速度$u$向右运动，那么这等价于装置不动，而以太风以速度$u$向左吹来。
容易得出光在水平方向从$O$到M$_1$再返回$O$的路径上所花费的时间是：
\[
	T_1 = \frac L{c - u} + \frac L{c + u} = \frac{2L/c}{1 - u^2/c^2}.
\]
同理，容易得出光在竖直方向上从$O$到M$_2$再返回$O$的路径上所花费的时间是：
\[
	T_2 = \frac{2L}{\sqrt{c^2 - u^2}} = \frac{2L/c}{\sqrt{1 - u^2/c^2}}.
\]
容易发现$T_1 \neq T_2$。由于装置在以太中的运动速度$u \ll c$，%对小量做Taylor展开得出光沿着两条路径上传播的时间差是：
\[
	T_1 - T_2 = \frac{2L}c\biggp{1 + \frac{u^2}{c^2}} - \frac{2L}c\biggp{1 + \frac{u^2}{2c^2}} + \bigo(u^4) = \frac{Lu^2}{c^3} + \bigo(u^4).
\]
将装置转过$\ang{90}$，两臂互换，模型预言两束光的相对时间延迟改变约$2Lu^2/c^3$，因而应有随装置方位变化的条纹位移。实验在其精度内未观察到这一模型预言的位移。

狭义相对论(special relativity)以相对性原理和真空光速不变为基本假设，为此类实验与电磁学提供了统一描述。

\section{狭义相对论}
\label{sec:special relativity}

\paragraph{Einstein的两个假设}

此节仅介绍必要内容，而略去狭义相对论的大部分推导内容。%如尺缩效应、钟慢效应等。

\begin{theorem}
    {相对性假设}{postulate of ralativity}
    物理定律在所有惯性系(inertial reference frame)中具有相同形式。惯性系之间作匀速直线相对运动，在一个封闭实验室中进行的实验不能确定其整体的绝对速度。
\end{theorem}

\begin{theorem}
    {光速不变假设}{postulate of the constancy of the speed of light}
    真空中的光在所有惯性系中均以相同的有限速率$c$传播，与光源和观察者的运动无关。
\end{theorem}

\subsection{Lorentz变换}

\begin{theorem}
    {Lorentz变换}{Lorentz transformation}
    给定惯性系$S$和$S'$，$S'$系相对$S$系有$x$方向的相对速度$v$。
    同一事件(event)在$S,S'$系中的坐标分别为$(ct,x,y,z)$和$(ct',x',y',z')$，则Lorentz变换给出：
    %\footnote{为了之后计算方便，也为了更好地凸显出理论的对称性，我们把$c$这个常数定义成1，也就是把光速$c$定义成一个标准单位．这样处理后时间和空间的量纲也统一起来了}
    \begin{equation}
    	\begin{bmatrix}
    		ct'\\x'\\y'\\z'
        \end{bmatrix} =
    	\begin{bmatrix}
        \gamma&-\gamma\beta\\-\gamma\beta&\gamma\\ &&1\\ &&&1
    	\end{bmatrix}
    	\begin{bmatrix}
    		ct\\x\\y\\z
    	\end{bmatrix}
    \end{equation}
    其中相对速度$\beta \equiv v/c$，Lorentz因子
    \begin{equation}
        \label{eq:ed-lorentz-factor}
        \gamma \equiv \frac1{\sqrt{1 - \beta^2}}.
    \end{equation}
\end{theorem}

如果质点在$S$系中有$x$方向的速度$u = \d x/\d t$，则其在$S'$中的速度为
\begin{equation}
    u' = \dv{x'}{t'} = \frac{\gamma(\d x - v\d t)}{\gamma(\d t - v\d x/c^2)} = \frac{u - v}{1 - vu/c^2}.
\end{equation}

\paragraph{四矢量}

为讨论形式的统一，我们定义四矢量(4-vector)。

\begin{definition}{四矢量}{4-vector}
    四矢量的逆变分量(contravariant components)不仅有四个分量，还必须按$a'{}^\mu = \Lambda^\mu{}_\nu a^\nu$变换，记为
    \begin{equation}
        a^\mu \equiv (a^0,a^1,a^2,a^3)\transp,
    \end{equation}
    % 称$a^0$是时间(temporal)分量。

    与之对应的是协变量(covariant)，形如
    \begin{equation}
        a_\mu \equiv (a_0,a_1,a_2,a_3).
    \end{equation}
\end{definition}
时空连续统(continuum)由位置四矢量表示：
\begin{equation}
    x^\mu \equiv (ct,x,y,z)\transp.
\end{equation}
\begin{definition}{Einstein求和约定}{Einstein summation convetion}
    表达式上下标中出现重复指标时，表示对其遍历求和，即
    \begin{equation}
        a_\mu b^\mu \equiv \sum_\mu a_\mu b^\mu,
    \end{equation}
    因此Einstein约定：这个求和符号是可以省略的。
\end{definition}
% 使用Einstein求和约定，定义四维标量积
% \[
%     a_\mu b^\mu:=a_0b^0+a_1b^1+a_2b^2+a_3b^3\equiv a^\mu b_\mu.
% \]
则Lorentz变换可以表示为
\begin{subequations}
    \label{eq:ed-boost-matrix}
    \begin{align}
        \bar x^\mu &= \Lambda^\mu{}_\nu x^\nu,\\
        \Lambda^\mu{}_\nu &\equiv
            \begin{bmatrix}
                \gamma&-\gamma\beta\\-\gamma\beta&\gamma\\ &&1\\ &&&1
            \end{bmatrix}
    \end{align}
\end{subequations}

\paragraph{时空结构}

狭义相对论时空的具体几何结构由间隔(interval)确定
\begin{equation}
    \label{eq:ed-invariant-interval}
    s^2 = -c^2(t_A - t_B)^2 + (x_A - x_B)^2 + (y_A - y_B)^2 + (z_A - z_B)^2.
\end{equation}
可以验证：间隔在Lorentz变换下是不变的：
\footnote{
不变量(invariant)表示在所有惯性系中均保持相同的值；守恒量(conserved quantity)表示在某个过程前后保持相同的值。比如，质量是不变量，却不守恒；能量是守恒的，但不是不变量；电荷量既是不变量，又是守恒量；而速度既不是不变量，也不守恒。

在每个封闭系统中，总的相对论能量和动量都是守恒的。
}
\[
    (s^2)' = -c^2(t_A' - t_B')^2 + (x_A' - x_B')^2 + (y_A' - y_B')^2 + (z_A' - z_B')^2 = s^2.
\]
\begin{itemize}
    \item 若$s^2 < 0$，则称间隔是类时的(timelike)，因为这是同地不同时事件间隔的符号；
    \item 若$s^2 > 0$，则是类空的(spacelike)，因为这是同时不同地事件间隔的符号；
    \item 若$s^2 = 0$，则是类光的(lightlike)，因为这是当两个事件由一个以光速传播的信号连接时所保持的关系。
\end{itemize}
\begin{definition}
    {度规张量}{metric tensor}
    在微分形式上，无穷小间隔是
    \begin{equation}
        \d s^2 = g_{\mu\nu}\d x^\mu\nd x^\nu.
    \end{equation}
    式中$g_{\mu\nu}$是度规张量(metric tensor)。对于狭义相对论的平坦时空（广义相对论中为弯曲时空），采取Minkowski度规：
    \begin{equation}
        g_{\mu\nu} \equiv
        \begin{bmatrix}
            -1\\ &1\\ &&1\\ &&&1
        \end{bmatrix}.
    \end{equation}
    逆变度规张量$g^{\mu\nu}$应满足
    \begin{equation}
        g_{\mu\nu}g^{\nu\rho} = \delta_\mu{}^\rho.
    \end{equation}
    % 为$g_{\mu\nu}$的归一化余子式。
    在这里所用的惯性系直角坐标中，度规矩阵等于其逆矩阵。
\end{definition}
\begin{remark}
    Minkowski度规有两种：
    \begin{itemize}
        \item 东海岸度规$\eta_{\mu\nu} = \diag(-1,1,1,1)$，常用于弦论和宇宙学，其好处是空间分量上的度规分量和我们平时习惯的Euclid空间的度规相同，Griffiths采用此；
        \item 西海岸度规$\eta_{\mu\nu} = \diag(1,-1,-1,-1)$，常用于粒子物理学，Jackson采用此；
        \item 其实还有几乎没人喜欢的$\i ct$。
    \end{itemize}
\end{remark}

\begin{corollary}
    协变坐标$x_\mu$可由逆变坐标$x^\mu$和度规张量得到
    \begin{equation}
        x_\mu = g_{\mu\nu}x^\nu.
    \end{equation}
    同样有$x^\mu = g^{\mu\nu}x_\nu$，写成分量形式就是
    \[
        x_0 = -x^0,\enspace x_1 = x^1,\enspace x_2 = x^2,\enspace x_3 = x^3.
    \]
\end{corollary}

\subsection{相对论动力学}

\paragraph{相对论运动学}

沿粒子的类时世界线，共动时钟测得的时间称为固有时间(proper time)。若粒子在惯性系中速度为$u$，则
\begin{equation}
    \label{eq:ed-proper-time}
    \d\tau = \sqrt{1 - u^2/c^2}\,\d t \equiv \frac{\d t}{\gamma_u}.
\end{equation}
进而定义四矢量速度
\begin{equation}
    \eta^\mu \equiv \dv{x^\mu}\tau,
\end{equation}
和四矢量动量
\begin{equation}
    p^\mu \equiv m\eta^\mu.
\end{equation}
其中$m$为(静止)质量。

$p^\mu$的时间分量是$p^0 = \gamma_umc$，对应相对论能量：
\begin{equation}
    \label{eq:ed-relativistic-energy}
    E = \gamma_umc^2,
\end{equation}
静止时，$\gamma_u = 1,\enspace E = mc^2$，二者之差定义为动能
\[
    \Ek \equiv E - mc^2 = \frac12mu^2 + \frac38\frac{mu^4}{c^2} + \frac5{16}\frac{mu^6}{c^4} + \bigo(u^8),
\]
相对论动量
\begin{equation}
    \label{eq:ed-relativistic-momentum}
    \bm p = m\bm\eta = m\dv{\bm r}\tau = \gamma_um\dv{\bm r}t = \gamma_um\bm u.
\end{equation}
有关系
\begin{equation}
    E^2 = p^2c^2 + m^2c^4.
\end{equation}
\begin{example}{Compton散射}{Compton scattering}
    频率为$\nu$的光子与一个静止的电子发生Compton散射，散射后光子的频率变为$\nu'$，散射角为$\theta$，反冲电子动量大小为$p_\electron$，反冲角为$\phi$。
    \begin{notecaptionblock}
        \tikzfigure{compton-scattering}
        \captionof{figure}{Compton散射示意图}
    \end{notecaptionblock}
    由横向、纵向的动量守恒以及能量守恒：
    \begin{subequations}
        \begin{align}
            \frac{h\nu}c &= \frac{h\nu'}c\cos\theta + p_\electron\cos\phi,\\
            0 &= \frac{h\nu'}c\sin\theta - p_\electron\sin\phi,\\
            h\nu + m_\electron c^2 &= h\nu' + \sqrt{(m_\electron c^2)^2 + (p_\electron c)^2}.
        \end{align}
    \end{subequations}
    消去$\phi$，
    解得
    \begin{equation}
        h\nu' = \frac{h\nu}{1 + \alpha(1 - \cos\theta)},\quad\alpha \equiv \frac{h\nu}{m_\electron c^2}.
    \end{equation}
    将上式改写成波长的形式，得到Compton移动
    \begin{equation}
        \Delta\lambda = \frac h{m_\electron c}(1 - \cos\theta).
    \end{equation}
    这个结果给定角度后就确定了光子能量。散射到该角度的概率由\secref{sec:qed Compton}的量子振幅给出。
\end{example}

\paragraph{相对论动力学}

相对论中的力依然遵循Newton第二定律中的定义：
\[
    \bm F \equiv \dv{\bm p}t,
\]
因此功能关系仍然满足：
\begin{equation}
    W \equiv \int\bm F \cdot \d\bm\ell = \int\dv{\bm p}t \cdot \bm u\d t = \int\dd t\frac{m\bm u}{\sqrt{1 - u^2/c^2}} \cdot \bm u\d t = \Delta E.
\end{equation}
而Newton第三定律在不同地的时候不再成立，因为在相对论中，同时性是相对的。

可得力在Lorentz变换中的关系
\begin{subequations}
    \begin{align}
        F_x' &= \frac{F_x - v\bm u \cdot \bm F/c^2}{1 - vu_x/c^2},\\
        F_y' &= \frac{F_y}{\gamma(1 - vu_x/c^2)},\\
        F_z' &= \frac{F_z}{\gamma(1 - vu_x/c^2)}.
    \end{align}
\end{subequations}
如果质点在$S$系中保持静止($u = 0$)，则其受到的力的变换为：
\begin{equation}
    \bm F_\perp' = \frac1\gamma\bm F_\perp,\qquad F_\parallel' = F_\parallel.
\end{equation}
\subsectionstar{Thomas进动}

粒子沿曲线运动时，它在相邻两点的瞬时静止系之间要作不同方向的Lorentz变换。两次不共线的纯推进合成后带有一个空间转动。一个在自身静止系中不受力矩的陀螺，由实验室逐时比较其指向，便可看出这个转动。

令$\bm s$为用实验室系沿瞬时速度$\bm u$作一次纯推进后测得的静止系自旋，四矢量$S^\mu$在该系中为$(0,\bm s)$。记$\gamma = \gamma_u$，变回实验室有
\begin{align*}
    S^0 &= \frac{\gamma}{c}\bm u \cdot \bm s,\\
    \bm S &= \bm s + \frac{\gamma^2}{(\gamma + 1)c^2}(\bm u \cdot \bm s)\bm u.
\end{align*}
由此$S^\mu\eta_\mu = 0$。无力矩意味着在瞬时静止系中$\d S^\mu/\nd\tau$的空间分量为零，故此导数平行于四维速度$\eta^\mu$。对正交关系求导，用$\eta^\mu\eta_\mu = -c^2$决定比例系数：
\begin{equation}
    \label{eq:ed-fermi-walker-transport}
    \dv{S^\mu}{\tau} = \frac{S_\nu a^\nu}{c^2}\eta^\mu,\qquad a^\mu = \dv{\eta^\mu}{\tau}.
\end{equation}
这是无力矩自旋的Fermi--Walker输运。用$\bm s = \bm S - \gamma S^0\bm u/[(\gamma + 1)c]$消去实验室四矢量的分量，并用$\dot\gamma = \gamma^3\bm u \cdot \bm a/c^2$，得到
\begin{subequations}
    \begin{align}
        \dv{\bm s}t &= \bm\omega_\mathrm{T} \times \bm s,\\
        \bm\omega_\mathrm{T} &= \frac{\gamma^2}{\gamma + 1}\frac{\bm a \times \bm u}{c^2},
        \quad \bm a = \dv{\bm u}t .
    \end{align}
\end{subequations}
匀速圆周运动中，取轨道角速度为$\bm\Omega$，则$\bm a \times \bm u = -u^2\bm\Omega$，从而$\bm\omega_\mathrm{T} = -(\gamma - 1)\bm\Omega$。每绕行一圈，陀螺相对于实验室反向进动$2\pi(\gamma - 1)$。实际带磁矩的粒子在磁场中还会受到磁力矩。

\section{相对论电动力学}
为什么会有磁场？利用静电学和相对论，我们可以计算出载电流线和运动电荷之间的磁力，而不需要援引磁力定律。
\begin{example}{导线的磁场力}{}
    $S$系中有一导线，其中正负电荷静止时的线密度为$\pm\lambda_0$，通电正负电荷后在导线中反向运动，速度$v_\pm = \pm v$，则正负电荷的线密度为$\lambda_\pm = \pm\gamma\lambda_0$，记$\lambda = \gamma\lambda_0 > 0$，电流为$I = 2\lambda v$，导线外距离$r$处有一电荷$q$沿平行于导线的方向以速度$u < v$运动。

    %$S$系中导线没有净电荷，因此电荷$q$不受电场力。
    \begin{notecaptionblock}
        \tikzfigure{q-in-lab-frame}
        \tikzfigure{q-in-q-frame}
        \captionof{figure}{$S$系和$S'$系中导线与电荷示意图}
    \end{notecaptionblock}
    转换到粒子静止的参考系$S'$，则电荷$q$不受磁场力，仅受电场力。
    $S'$系中正负电荷的速度为
    \[
        v_\pm' = \frac{v_\pm - u}{1 - v_\pm u/c^2} = \frac{\pm v - u}{1 \mp vu/c^2},
    \]
    正负电荷线密度变为$\lambda_\pm' = \pm\gamma_{v_\pm'}\lambda_0$，总电荷线密度和电场为
    \begin{align*}
        \lambda_\tot &= -\frac{2vu}{c^2}\gamma_u\lambda,\\
        E &= \frac{\lambda_\tot}{2\pi\varepsilon_0r}.
    \end{align*}
    因此$S'$系中电荷所受的电场力为
    \[
        F' = qE = -\frac{vu\gamma_u\lambda}{c^2}\frac q{\pi\varepsilon_0r}.
    \]
    变回到$S$系中，
    \[
        F = \frac1{\gamma_u}F' = -\frac{vu\lambda}{c^2}\frac q{\pi\varepsilon_0r} = -qu\frac{\mu_0I}{2\pi r}.
    \]
    而$S$系中导线没有净电荷，$q$不受电场力，故此力只能是磁场力！
\end{example}
\paragraph{场的变换}
假设：
\begin{compactenum}
	\item 电荷是不变的；
	\item 无论场是如何产生的，转换规则都是相同的。
\end{compactenum}
得到
\begin{align}
    E_x' &= E_x,&E_y' &= \gamma(E_y - vB_z),&E_z' &= \gamma(E_z + vB_y);\\
    B_x' &= B_x,&B_y' &= \gamma\Bigp{B_y + \frac v{c^2}E_z},&B_z' &= \gamma\Bigp{B_z - \frac v{c^2}E_y}.
\end{align}
\begin{example}{匀速运动电荷的场}{E, B of q with v}
    考虑$S'$系中，电荷是静止的，没有磁场只有电场：
    \[
        E' = \frac1{4\pi\varepsilon_0}\frac q{r'^2}\uvec r',
    \]
    变换到原系$S$：
    \[
        E_x = E_x',\quad E_y = \gamma E_y',\quad E_z = \gamma E_z',
    \]
    得到电场为$\uvec r$方向，但并不各向同性：
    \begin{equation}
        \bm E = \frac1{4\pi\varepsilon_0}\frac{1 - \beta^2}{(1 - \beta^2\sin^2\theta)^{3/2}}\frac q{r^2}\uvec r.
    \end{equation}
    磁场为$\uvec\phi$方向：
    \begin{equation}
        \bm B = \frac{\bm v}{c^2} \times \bm E = \frac{\mu_0}{4\pi}\frac{(1 - \beta^2)\sin\theta}{(1 - \beta^2\sin^2\theta)^{3/2}}\frac{qv}{r^2}\uvec\phi.
    \end{equation}
\end{example}
\subsection{电磁场张量}
电磁场的变换是由一个反对称的二阶张量连接起来的
\[
    t^{\mu\nu} =
    \begin{bmatrix}
        0&t^{01}&t^{02}&t^{03}\\
        -t^{01}&0&t^{12}&t^{13}\\
        -t^{02}&-t^{12}&0&t^{23}\\
        -t^{03}&-t^{13}&-t^{23}&0
    \end{bmatrix}.
\]
变换为 
\[
    \bar t^{\mu\nu} = \Lambda^\mu{}_\rho\Lambda^\nu{}_\sigma t^{\rho\sigma}.
\]
对应分量的变换规则为
\begin{align*}
    \bar t^{01} &= t^{01},&\bar t^{02} &= \gamma(t^{02} - \beta t^{12}),&\bar t^{03} &= \gamma(t^{03} + \beta t^{31}),\\
    \bar t^{23} &= t^{23},&\bar t^{31} &= \gamma(t^{31} + \beta t^{03}),&\bar t^{12} &= \gamma(t^{12} - \beta t^{02}).
\end{align*}
与场变换比较，可得电磁张量
\begin{equation}
    F^{\mu\nu} =
    \begin{bmatrix}
        0&E_x/c&E_y/c&E_z/c\\
        -E_x/c&0&B_z&-B_y\\
        -E_y/c&-B_z&0&B_x\\
        -E_z/c&B_y&-B_x&0
    \end{bmatrix}.
\end{equation}
或对偶(dual)张量
\begin{equation}
    G^{\mu\nu} =
    \begin{bmatrix}
        0&B_x&B_y&B_z\\
        -B_x&0&-E_z/c&E_y/c\\
        -B_y&E_z/c&0&-E_x/c\\
        -B_z&-E_y/c&E_x/c&0
    \end{bmatrix}.
\end{equation}
\paragraph{张量符号中的电动力学}
用张量的语言重新表述电动力学定律（Maxwell方程和Lorentz力）。统一电荷源$\rho$和电流源$\bm J$，定义四矢量电流源
\begin{equation}
    J^\mu = (c\rho,J_x,J_y,J_z)\transp,
\end{equation}
\begin{definition}{四维导数}{}
    定义 
    \begin{equation}
        \p_\mu \equiv \pp{x^\mu},\quad\p^\mu \equiv \pp{x_\mu}.
    \end{equation}
    有
    \begin{equation}
        \p_\mu a^\mu = \p_0a^0 + \p_1a^1 + \p_2a^2 + \p_3a^3 = \pv{a^0}{x^0} + \pv{a^1}{x^1} + \pv{a^2}{x^2} + \pv{a^3}{x^3}.
    \end{equation}
\end{definition}
连续性方程\eqref{eqn:continuity}变为 
\begin{equation}
    \label{eq:ed-covariant-continuity}
    \p_\mu J^\mu = \pv\rho t + \div\bm J = 0,
\end{equation}
Maxwell方程组变为 
\begin{subequations}
    \begin{align}
        \p_\nu F^{\mu\nu} &= \mu_0J^\mu,\\
        \p_\nu G^{\mu\nu} &= 0.
    \end{align}
\end{subequations}
$\mu = 0$给出
\begin{subequations}
    \begin{align}
        \div\bm E &= \rho/\varepsilon_0,\\
        \div\bm B &= 0;
    \end{align}
\end{subequations}
$\mu = 1,2,3$给出$x,y,z$分量
\begin{subequations}
    \begin{align}
        \curl\bm B - \mu_0\varepsilon_0\pv{\bm E}t &= \mu_0\bm J,\\
        \curl\bm E + \pv{\bm B}t &= \bm0.
    \end{align}
\end{subequations}

定义四矢量Minkowski力 
\begin{equation}
    K^\mu := q\eta_\nu F^{\mu\nu},
\end{equation}
得到Lorentz力的表达式
\[
    \bm F = q(\bm E + \bm u \times \bm B),
\]

\subsection{相对论势}

将标量势和矢量势的定义：
\begin{align*}
    \bm E &= -\grad\Phi - \pv{\bm A}t,\\
    \bm B &= \curl\bm A,
\end{align*}
展开，可见对称性：
\begin{alignat*}{2}
    E_x &= -\pv\Phi x - \pv{A_x}t,&\qquad B_x &= \pv{A_z}y - \pv{A_y}z,\\
    E_y &= -\pv\Phi y - \pv{A_y}t,&B_y &= \pv{A_x}z - \pv{A_z}x,\\
    E_z &= -\pv\Phi z - \pv{A_z}t,&B_z &= \pv{A_y}x - \pv{A_x}y,
\end{alignat*}

四矢量势
\begin{equation}
    A^\mu := (\Phi/c,A_x,A_y,A_z)\transp,
\end{equation}
则电磁张量可以写成
\begin{equation}
    F^{\mu\nu} = \p^\mu A^\nu - \p^\nu A^\mu.
\end{equation}
% 注意这里是协变量$x_\mu$，$x_0=-ct$有一个负号。
\begin{example}{匀速运动电荷的场·续}{E, B of q with v, II}
    续\exmref{exm:E, B of q with v}，我们用势的方法再做一遍。令两坐标系的原点在$t = t' = 0$重合，$S'$系中电荷静止于原点
    \[
        \Phi' = \frac1{4\pi\varepsilon_0}\frac q{r'},\quad\bm A' = \bm 0,
    \]
    变回$S$系
    \[
        \Phi = \gamma\Phi',\quad A_x = \gamma\frac v{c^2}\Phi',\quad A_y = A_z = 0,
    \]
    其中 
    \[
        r' = \sqrt{x'^2 + y'^2 + z'^2} = \sqrt{\gamma^2(x - vt)^2 + y^2 + z^2}.
    \]
    则
    \[
        \bm E = -\grad\Phi - \pv{\bm A}t = \frac1{4\pi\varepsilon_0}\frac{\gamma q\bm R}{\bigb{\gamma^2(x - vt)^2 + y^2 + z^2}^{3/2}}.
    \]
    其中$\bm R = (x - vt,y,z)$。
\end{example}
利用四矢量势，Maxwell方程组变为
\[
    \p^\mu\p_\nu A^\nu - \p^\nu\p_\nu A^\mu = \mu_0J^\mu,
\]
使用Lorenz规范
\begin{equation}
    \label{eq:ed-covariant-lorenz-gauge}
    \p_\mu A^\mu = \frac1{c^2}\pv\Phi t + \div\bm A = 0.
\end{equation}
由此我们得到了Maxwell方程组最简单的形式：
\begin{equation}
    \p^\nu\p_\nu A^\mu = -\mu_0J^\mu.
\end{equation}

\section{外电磁场中带电粒子的运动}
设质量为$m$、电荷为$q$的粒子在给定的外场中运动，忽略粒子自身辐射对轨道的影响。沿轨迹的作用量为
\begin{subequations}
    \label{eq:ed-particle-action}
    \begin{align}
        S &= -mc^2\int\d\tau + q\int A_\mu\,\d x^\mu
        = \int L\,\d t,\\
        L &= -mc^2\sqrt{1 - u^2/c^2} + q\bm u \cdot \bm A - q\Phi .
    \end{align}
\end{subequations}
因为$A_0 = -\Phi/c$，静止粒子的电势能为$q\Phi$。正则动量$\bm P = \p L/\p\bm u$和Hamilton量$H = \bm P \cdot \bm u - L$分别是
\begin{subequations}
    \begin{align}
        \bm P &= \gamma_um\bm u + q\bm A,\\
        H &= \sqrt{m^2c^4 + c^2\abs{\bm P - q\bm A}^2} + q\Phi.
    \end{align}
\end{subequations}
这里$\bm p = \bm P - q\bm A$就是前面由速度定义的机械动量。势作规范变换$\bm A \mapsto \bm A + \grad\chi,\ \Phi \mapsto \Phi - \p_t\chi$时，$L$增加沿轨迹的全导数$q\,\d\chi/\nd t$，固定端点的变分不变。写出Euler--Lagrange方程的第$i$个分量，
\[
    \dv{p_i}t = q\biggl[-\p_i\Phi - \p_t A_i
               + \sum_j u_j(\p_i A_j - \p_j A_i)\biggr],
\]
括号中最后一项正是$(\bm u \times \bm B)_i$，从作用量便得到Lorentz力。

\begin{example}{正交均匀电磁场中的漂移}{crossed field drift}
    设$\bm E \perp \bm B$且$E < cB$。取一个以$\bm v_\mathrm{d} = \bm E \times \bm B/B^2$运动的惯性系。由于$\bm v_\mathrm{d} \times \bm B = -\bm E$，由场变换可知$\bm E' = 0$，且$B' = \sqrt{B^2 - E^2/c^2}$。在这个系中，磁力保持粒子能量不变，沿磁场的动量也守恒。垂直磁场的动量匀速转动，轨道半径和回旋角频率分别为
    \[
        r_\mathrm{L}' = \frac{p_\perp'}{\abs q B'},\qquad
        \omega_\mathrm{B}' = \frac{\abs q B'}{\gamma' m}.
    \]
    一个回旋周期后，粒子在该系的横向位移为零。变回实验室，横向平均速度即为$\bm v_\mathrm{d}$，与电荷正负无关。同一观察也给出条件$E < cB$的物理意义：只有这时，使电场消失的参考系所需速率$E/B$才小于$c$。
\end{example}

\subsection{电磁场的作用量}
上面固定势，对粒子轨迹作变分。反过来，给定满足连续性方程的四维电流$J^\mu$，也可以把势$A_\mu$作为待求的场。在SI中，真空电磁场及其与源的耦合对应的作用量为
\begin{equation}
    \label{eq:ed-electromagnetic-action}
    S[A] = \int\biggl(-\frac{1}{4\mu_0}F_{\mu\nu}F^{\mu\nu} + J^\mu A_\mu\biggr)\d V\,\d t .
\end{equation}
第一项等于$\varepsilon_0 E^2/2 - B^2/(2\mu_0)$，第二项等于$\bm J \cdot \bm A - \rho\Phi$。由于$F^{\mu\nu}$反对称，在边界上固定势并分部积分，得到
\[
    \vd S = \int\biggl(\frac{1}{\mu_0}\p_\mu F^{\mu\nu} + J^\nu\biggr)\vd A_\nu\,\d V\,\d t .
\]
要求$\vd S = 0$即得有源Maxwell方程$\p_\mu F^{\nu\mu} = \mu_0 J^\nu$。另两条无源方程由$F^{\mu\nu} = \p^\mu A^\nu - \p^\nu A^\mu$的定义自动满足。

第\pageref{par:photon mass test}页的球壳实验可以从同一作用量理解。Proca理论在场的Lagrange密度中加入一项
\[
    \Delta\mathcal L = -\frac{\kappa^2}{2\mu_0}A_\mu A^\mu .
\]
对势作变分，场方程随之变为
\begin{equation}
    \label{eq:ed-proca-field}
    \p_\mu F^{\nu\mu} + \kappa^2 A^\nu = \mu_0J^\nu .
\end{equation}
取四维散度，反对称项相消，由电荷守恒得$\kappa^2\p_\nu A^\nu = 0$。当$\kappa \ne 0$时，$\p_\nu A^\nu = 0$由动力学强制满足。静电极限的时间分量给出$(\laplacian - \kappa^2)\Phi = -\rho/\varepsilon_0$，正是球壳内部电势偏离常量所依据的方程。

\subsection{电磁场的能动张量与辐射压力}
\label{sec:field stress energy}
一束光照在镜面上，反射光带走能量，镜面却仍然受力。能动张量(energy-momentum tensor)把能量和三分量动量的输运写在同一个局域守恒律中。现在回到无质量光子的Maxwell场，取$x^0 = ct$、$g_{\mu\nu} = \diag(-1,1,1,1)$。

可以由场的作用量直接求得对称能动张量。把度规看作计算用的变量，记$g = \det(g_{\mu\nu})$、$\mathcal L_\mathrm{f} = -F_{\alpha\beta}F^{\alpha\beta}/(4\mu_0)$。对$\sqrt{-g}\mathcal L_\mathrm{f}$求变分时，固定$A_\mu$，故$F_{\alpha\beta}$固定，抬指标的度规参与变分。利用$\vd\sqrt{-g} = -\sqrt{-g}\,g_{\mu\nu}\vd g^{\mu\nu}/2$，得
\begin{equation}
    T^{\mu\nu}_\mathrm{em}
    = \frac{1}{\mu_0}\biggl(F^{\mu\lambda}F^\nu{}_\lambda
       - \frac14 g^{\mu\nu}F_{\alpha\beta}F^{\alpha\beta}\biggr),
    \label{eqn:field stress energy}
\end{equation}
其中能动张量的符号由变分关系确定，
\begin{equation}
    \label{eq:ed-metric-stress-definition}
    \vd(\sqrt{-g}\mathcal L_\mathrm{f})
      = -\frac{\sqrt{-g}}{2}\,T_{\mu\nu}^\mathrm{em}\vd g^{\mu\nu}.
\end{equation}
将平直时空的场张量代回，$F_{\alpha\beta}F^{\alpha\beta} = 2(B^2 - E^2/c^2)$，于是
\begin{subequations}\label{eqn:field stress components}
    \begin{align}
        T^{00}_\mathrm{em} &= \frac{\varepsilon_0 E^2}{2} + \frac{B^2}{2\mu_0} = u,\\
        T^{0i}_\mathrm{em} = T^{i0}_\mathrm{em} &= \frac{S_i}{c} = c\,g_i,\\
        T^{ij}_\mathrm{em} &= u\delta_{ij} - \varepsilon_0 E_iE_j - \frac{B_iB_j}{\mu_0} = -T_{ij}^{\mathrm{Maxwell}}.
    \end{align}
\end{subequations}
分量$T^{ij}_\mathrm{em}$给出沿$i$方向输运的$j$分量动量通量。它等于式\eqref{eqn:vacuum maxwell stress}中Maxwell应力的负值。场的局域守恒律为
\begin{equation}
    \label{eq:ed-covariant-stress-balance}
    \p_\mu T^{\mu\nu}_\mathrm{em} = -F^{\nu\lambda}J_\lambda .
\end{equation}
取$\nu = 0$并乘$c$即得$\p_t u + \div\bm S = -\bm J \cdot \bm E$。取$\nu = i$则得$\p_t g_i + \p_j T^{ji}_\mathrm{em} = -(\rho\bm E + \bm J \times \bm B)_i$，右端是场传给物质的力密度的负值。

\begin{example}{平面镜所受的辐射压力}{pressure on a mirror}
    在真空中令线偏振平面波沿$+z$正入射到$z = 0$的静止理想导体，入射电场峰值为$E_0$。在镜面，入射和反射电场相消，两波的切向磁场同向，故$E = 0$，$B_y = 2(E_0/c)\cos\omega t$。导体内的振荡场为零。由式\eqref{eqn:field stress components}，单位时间进入镜面的法向动量除以面积，给出时均压力
    \begin{equation}
        \label{eq:ed-perfectly-reflecting-radiation-pressure}
        p_\mathrm{rad} = \ave{T^{zz}_\mathrm{em}(0^-)}
         = \frac{\ave{B_y^2}}{2\mu_0}
         = \varepsilon_0 E_0^2 = \frac{2I}{c}.
    \end{equation}
    入射波的平均强度为$I = \varepsilon_0 c E_0^2/2$。
    镜面处的时均净能流为零。入射能量$U$携带的法向动量为$U/c$，反射后变为$-U/c$，镜面正好获得$2U/c$。
\end{example}
